% =====================================================================
% GLOSSAIRE DES SIGLES
%
% PIÈCE PROPRE À LA VERSION ENSAE, placée en tête avant les notes de synthèse, sur
% le modèle du mémoire de F. Dountio (ENSAE, 2026). La table des notations
% mathématiques reste en annexe : ce glossaire ne porte que les sigles et les
% termes qu'un lecteur non spécialiste rencontre dès les notes de synthèse.
% =====================================================================

\chapter*{Glossaire}
\addcontentsline{toc}{chapter}{Glossaire}
\markboth{Glossaire}{}

% HARNAIS-HORS-SECTION: glossaire de sigles ; les seuls nombres sont des numeros
% de pilier, un numero de reglement et un niveau de quantile reglementaire.

\begin{center}\small
\renewcommand{\arraystretch}{1.15}
\begin{longtable}{@{}>{\raggedright\arraybackslash}p{2.6cm}>{\raggedright\arraybackslash}p{12.2cm}@{}}
\toprule
\textbf{Sigle ou terme} & \textbf{Signification} \\
\midrule
\endhead
ACPR & Autorité de contrôle prudentiel et de résolution \\
AMRAE & Association pour le management des risques et des assurances de l'entreprise, qui publie l'étude LUCY \\
BSCR & \emph{Basic Solvency Capital Requirement}, capital de solvabilité de base de la Formule Standard \\
CTE & \emph{Conditional Tail Expectation}, espérance conditionnelle de la perte au-delà d'un quantile \\
DORA & \emph{Digital Operational Resilience Act}, règlement (UE) 2022/2554 sur la résilience opérationnelle numérique du secteur financier \\
EIOPA & Autorité européenne des assurances et des pensions professionnelles \\
Effet de fermeture & Capital retiré en corrigeant un seul canal à partir de l'état non conforme \\
Effet isolé & Capital ajouté par un seul canal à partir de l'état conforme \\
FFT & \emph{Fast Fourier Transform}, transformée de Fourier rapide, utilisée pour inverser la loi de la charge annuelle \\
GPD & \emph{Generalized Pareto Distribution}, loi de Pareto généralisée, loi des excès au-dessus d'un seuil élevé \\
Hackmageddon & Recensement public d'incidents cyber, utilisé pour la répartition par vecteur d'attaque \\
IC & Intervalle de confiance \\
Identification partielle & Démarche qui, faute de pouvoir estimer un paramètre, calcule l'ensemble de ses valeurs compatibles avec la donnée et publie des bornes \\
LUCY & \emph{Lumière sur la cyberassurance}, étude annuelle du marché français de la cyberassurance \\
ORSA & \emph{Own Risk and Solvency Assessment}, évaluation interne des risques et de la solvabilité \\
P1 à P5 & Les cinq piliers de DORA : gouvernance du risque TIC, gestion des incidents, tests de résilience, risque lié aux prestataires tiers, partage d'informations \\
PIT & \emph{Probability Integral Transform}, transformée intégrale de probabilité, utilisée pour tester une loi prédictive \\
POT & \emph{Peaks Over Threshold}, méthode d'estimation de queue par les dépassements d'un seuil \\
PRC & \emph{Privacy Rights Clearinghouse}, base publique de notifications de violations de données \\
SAS OpRisk Global & Base commerciale de pertes opérationnelles publiques, en montants \\
SCR & \emph{Solvency Capital Requirement}, capital de solvabilité requis sous Solvabilité~II, quantile à 99,5\,\% de la perte annuelle \\
SFCR & \emph{Solvency and Financial Condition Report}, rapport public sur la solvabilité et la situation financière \\
Shapley & Règle d'attribution qui partage un écart entre plusieurs causes en moyennant leurs contributions sur tous les ordres d'introduction \\
SLA & \emph{Single Loss Approximation}, approximation du quantile agrégé par celui d'une perte unique \\
TIC & Technologies de l'information et de la communication \\
TVaR & \emph{Tail Value-at-Risk}, moyenne de la perte au-delà de la VaR \\
VaR & \emph{Value-at-Risk}, quantile de la loi de perte ; dans ce mémoire, réservée à la charge annuelle agrégée \\
VERIS & \emph{Vocabulary for Event Recording and Incident Sharing}, schéma de description d'incidents ; VCDB en est la base communautaire \\
\bottomrule
\end{longtable}
\end{center}
